% Unit 102 terjemahan Bahasa Indonesia dari chapter7.tex baris 2472--2678.
% Sumber: Methods of Algebra, Volume 2, commit
% 9a5803ff77dd3257484cb177f851a73770a59dd3 (CC BY 4.0).
% stable-unit-id: o014.aljabr2.chapter7.galois-descent-h1
% Cakupan: Bagian 7.11 lengkap, hubungan penurunan Galois dengan kohomologi
% nonabelian H^1, hingga tepat sebelum Latihan.

% segment-id: o014.aljabr2.chapter7.galois-descent-h1.h001
\section{Kaitan antara Penurunan Galois dan
\texorpdfstring{$\Hm^1$}{H1}}\label{sec:Galois-H1}
\hypertarget{o014.aljabr2.chapter7.galois-descent-h1}{ }
% segment-id: o014.aljabr2.chapter7.galois-descent-h1.p001
Metode penurunan merupakan alat yang ampuh untuk mengklasifikasikan berbagai
objek aljabar atau geometri; pernyataan umumnya memerlukan bahasa geometri
aljabar. Demi kemudahan penyajian, dalam bagian ini kita membuat dua asumsi
penyederhanaan.
% segment-id: o014.aljabr2.chapter7.galois-descent-h1.l001
\begin{enumerate}
	\item Kita hanya menggunakan penurunan Galois
	(Teorema~\ref{prop:Galois-descent} beserta pembahasan sesudahnya), tanpa
	menggunakan penurunan datar atau teknik yang lebih umum.
	\item Buku ini sedapat mungkin memakai bahasa aljabar, khususnya aljabar
	linear, untuk menjelaskan contoh. Karena itu, ruang lingkup masalah
	klasifikasi harus dipersempit. Sesungguhnya, yang akan dibahas hanyalah
	satu kelas tensor dalam ruang vektor berdimensi hingga.
\end{enumerate}

% segment-id: o014.aljabr2.chapter7.galois-descent-h1.p002
Sudut pandang ini secara alami mengarah ke kohomologi nonabelian
$\Hm^1$ dari grup Galois. Pembahasan berikut bertumpu pada perangkat dari
\S\ref{sec:nonabelian-coh} dan \S\ref{sec:descent-Galois}. Karena kita perlu
mengizinkan perluasan Galois takhingga, kita juga tetap menggunakan bahasa
grup profinit; lihat \S\ref{sec:profinite-groups}. Pertama-tama, kita
perkenalkan notasi yang diperlukan.

% segment-id: o014.aljabr2.chapter7.galois-descent-h1.cv001
\begin{convention}
	Dalam bagian ini, tetapkan medan $K$ dan tuliskan
	$\otimes:=\otimes_K$. Untuk ruang vektor-$K$ berdimensi hingga $V$,
	tuliskan $V^\vee$ bagi ruang dualnya. Untuk setiap
	$(p,q)\in\Z_{\geq0}^2$, perkenalkan notasi
	% segment-id: o014.aljabr2.chapter7.galois-descent-h1.q001
	\[
	T^{p,q}V:=V^{\otimes p}\otimes(V^\vee)^{\otimes q}.
	\]
	Mengikuti kebiasaan dalam geometri diferensial, unsur-unsur
	$T^{p,q}V$ disebut \emph{tensor bertipe $(p,q)$ pada $V$}.
	\index{tensor}
\end{convention}

% segment-id: o014.aljabr2.chapter7.galois-descent-h1.p003
Setiap isomorfisme $g:V\to V'$ antara ruang vektor-$K$ berdimensi hingga
secara alami menginduksi
$T^{p,q}V\rightiso T^{p,q}V'$; demi kemudahan, isomorfisme yang diinduksi ini
tetap dilambangkan dengan $g$. Selain itu, sifat berdimensi hingga memberikan
isomorfisme kanonik
$(V_1\otimes V_2)^\vee\simeq V_1^\vee\otimes V_2^\vee$.

% segment-id: o014.aljabr2.chapter7.galois-descent-h1.p004
Dalam masalah klasifikasi yang dibahas pada bagian ini, objek utamanya adalah
data berbentuk berikut, yang didefinisikan di atas medan $K$:
% segment-id: o014.aljabr2.chapter7.galois-descent-h1.q002
\begin{equation}\label{eqn:tensor-data}
	\mathbf{t}=(V,I,(t_i,p_i,q_i)_{i\in I}),
\end{equation}
dengan
% segment-id: o014.aljabr2.chapter7.galois-descent-h1.l002
\begin{itemize}
	\item $V$: ruang vektor-$K$ berdimensi hingga;
	\item $I$: suatu himpunan (yang boleh kosong);
	\item $(p_i,q_i)$: keluarga pasangan bilangan bulat taknegatif, dengan
	$i$ merentang $I$;
	\item $t_i\in T^{p_i,q_i}V$: keluarga tensor, dengan $i$ merentang $I$.
\end{itemize}

% segment-id: o014.aljabr2.chapter7.galois-descent-h1.p005
Dengan kata lain, kita mempelajari suatu keluarga tensor bertipe tertentu
dalam ruang vektor berdimensi hingga. Untuk data lain
% segment-id: o014.aljabr2.chapter7.galois-descent-h1.q003
\[
\mathbf{t}'=(V',I,(t'_i,p_i,q_i)_{i\in I}),
\]
jika terdapat isomorfisme ruang vektor-$K$ $g:V\rightiso V'$ sedemikian
sehingga $gt_i=t'_i$ untuk setiap $i$, maka $g$ disebut isomorfisme antara
kedua data tersebut dan kita tuliskan
% segment-id: o014.aljabr2.chapter7.galois-descent-h1.q004
\[
g\mathbf{t}=\mathbf{t}'\quad\text{atau}\quad
g:\mathbf{t}\rightiso\mathbf{t}'.
\]

% segment-id: o014.aljabr2.chapter7.galois-descent-h1.p006
Kita hendak mengklasifikasikan data $\mathbf{t}$ ini hingga isomorfisme,
atau lebih lanjut mengklasifikasikan data $\mathbf{t}$ yang memenuhi
syarat-syarat tambahan. Syarat tersebut harus dapat dirumuskan secara
aljabar dan dipertahankan oleh isomorfisme. Meskipun kerangka teorinya
terbatas pada aljabar linear, jangkauan penerapannya sudah cukup luas.
Berikut beberapa kasus khusus.
% segment-id: o014.aljabr2.chapter7.galois-descent-h1.l003
\begin{enumerate}
	\item Tinjau data $(V,b)$, dengan
	$b\in(V^\vee)^{\otimes2}$ tensor bertipe $(0,2)$. Mengklasifikasikan data
	ini sama dengan mengklasifikasikan bentuk bilinear pada $V$. Kita juga
	dapat memberlakukan syarat bahwa $b$ simetris, antisimetris, takdegenerat,
	dan sebagainya; semuanya merupakan syarat ``aljabar''. Bentuk multilinear
	dapat diperlakukan secara serupa.
	\item Kita juga dapat mengklasifikasikan
	$(V,\varphi_1,\ldots,\varphi_n)$, dengan
	$\varphi_i\in V\otimes V^\vee$ tensor bertipe $(1,1)$. Karena
	$V\otimes V^\vee\simeq\End_K(V)$, ini sama dengan mengklasifikasikan $V$
	bersama suatu keluarga endomorfisme linear. Berbagai relasi aljabar dapat
	diberlakukan, misalnya komutativitas antara endomorfisme tersebut.
	\item Selanjutnya, tinjau data $(V,\mu)$, dengan
	$\mu\in V\otimes(V^\vee)^{\otimes2}$ tensor bertipe $(1,2)$, yang juga
	dapat dipandang sebagai unsur $\Hom_K(V\otimes V,V)$. Ini sama dengan
	mengklasifikasikan $V$ bersama suatu operasi ``perkalian'' bilinear
	$V\times V\to V$. Jika operasi itu disyaratkan asosiatif dan mempunyai
	unsur satuan (yang terakhir dapat dipandang sebagai tensor bertipe
	$(1,0)$), maka syarat tersebut sama dengan menjadikan $V$ aljabar-$K$.
	Kita juga dapat memberlakukan syarat lain pada $\mu$ untuk mempelajari
	beragam struktur seperti aljabar Lie dan aljabar Jordan.
\end{enumerate}

% segment-id: o014.aljabr2.chapter7.galois-descent-h1.p007
Secara lebih umum, \eqref{eqn:tensor-data} dapat diperluas sewajarnya
untuk membahas klasifikasi suatu keluarga subruang vektor, atau setidaknya
garis, di dalam $T^{p,q}V$. Karena keterbatasan ruang, hal ini tidak dibahas
lebih lanjut.

% segment-id: o014.aljabr2.chapter7.galois-descent-h1.d001
\begin{definition}
	Untuk ruang vektor-$K$ berdimensi hingga $V$, tuliskan $\GL(V)$ bagi grup
	automorfisme linearnya. Untuk data $\mathbf{t}$ seperti dalam
	\eqref{eqn:tensor-data}, definisikan subgrup dari $\GL(V)$
	% segment-id: o014.aljabr2.chapter7.galois-descent-h1.q005
	\[
	G(\mathbf{t}):=\{g\in\GL(V):g\mathbf{t}=\mathbf{t}\}.
	\]
\end{definition}

% segment-id: o014.aljabr2.chapter7.galois-descent-h1.p008
Pilih suatu basis bagi $V$ dan tuliskan syarat $gt_i=t_i$ dalam koordinat
untuk setiap $i\in I$. Dengan demikian, $G(\mathbf{t})$ dipotong di dalam
$\GL(V)$ oleh suatu keluarga persamaan polinomial. Dalam contoh-contoh di
atas, jika $\mathbf{t}$ bersesuaian dengan bentuk bilinear simetris (atau
antisimetris) takdegenerat $b:V\times V\to K$, maka $G(\mathbf{t})$ yang
bersesuaian ialah grup ortogonal $\mathrm{O}(V,b)$ (atau grup simplektik
$\mathrm{Sp}(V,b)$). Jika $\mathbf{t}$ adalah tensor bertipe
$(0,\dim V)$ yang diberikan oleh determinan, maka $G(\mathbf{t})$ yang
bersesuaian ialah grup linear khusus
$\SL(V):=\{g:\det g=1\}$.

% segment-id: o014.aljabr2.chapter7.galois-descent-h1.d002
\begin{definition}\label{def:tL}
	Untuk setiap perluasan medan $L|K$ dan ruang vektor-$K$ $V$, tuliskan
	$V_L:=L\dotimes{K}V$, yang dipandang sebagai ruang vektor-$L$. Melalui
	isomorfisme kanonik ruang vektor-$L$
	% segment-id: o014.aljabr2.chapter7.galois-descent-h1.q006
	\[
	(T^{p,q}V)_L\simeq T^{p,q}(V_L),
	\]
	setiap data $\mathbf{t}$ pada $V$ yang berbentuk
	\eqref{eqn:tensor-data} menginduksi data $\mathbf{t}_L$ pada $V_L$, yang
	didefinisikan di atas medan perluasan $L$.
\end{definition}

% segment-id: o014.aljabr2.chapter7.galois-descent-h1.p009
Perluasan medan tidak mengubah bagian numerik data tersebut,
$(\dim V,I,(p_i,q_i)_i)$.

% segment-id: o014.aljabr2.chapter7.galois-descent-h1.p010
Ingat bahwa bila $L|K$ merupakan perluasan Galois,
$\Gal(L|K)$ beraksi secara mulus dan semilinear pada $V_L$
(Lema~\ref{prop:Galois-descent-action}); aksi dituliskan dari kiri.
Sesungguhnya, grup itu beraksi secara mulus dan semilinear pada setiap
$T^{p,q}(V_L)$ (Definisi~\ref{def:semilinear-action}), secara ``diagonal''
pada hasil kali tensor dan melalui
$\check v\mapsto\check v\sigma^{-1}$ pada ruang dual $V^\vee$ (aksi
kontragradien, atau ``langkah invers''). Relatif terhadap isomorfisme di atas, subruang invariannya
adalah
% segment-id: o014.aljabr2.chapter7.galois-descent-h1.q007
\[
T^{p,q}(V_L)^{\Gal(L|K)}=T^{p,q}V.
\]
Ini tidak lain adalah penurunan Galois pada tingkat objek. Tingkat morfisme
dinyatakan oleh hasil berikut.

% segment-id: o014.aljabr2.chapter7.galois-descent-h1.lm001
\begin{lemma}
	Misalkan $V$ dan $W$ ruang vektor-$K$ berdimensi hingga. Untuk setiap
	$h\in\Hom_L(V_L,W_L)$ dan $\sigma\in\Gal(L|K)$, definisikan
	% segment-id: o014.aljabr2.chapter7.galois-descent-h1.q008
	\[
	{}^\sigma h:=\sigma h\sigma^{-1}:V_L\to W_L.
	\]
	Pemetaan ini tetap linear-$L$. Rumus tersebut memberikan aksi mulus dan
	semilinear $\Gal(L|K)$ pada ruang vektor-$L$
	$\Hom_L(V_L,W_L)$. Aksi ini kompatibel dengan komposisi pemetaan linear,
	dan
	% segment-id: o014.aljabr2.chapter7.galois-descent-h1.q009
	\[
	\Hom_L(V_L,W_L)^{\Gal(L|K)}=\Hom_K(V,W).
	\]
	Di sini $\Hom_K(V,W)$ ditanamkan ke dalam $\Hom_L(V_L,W_L)$ melalui
	$f\mapsto f_L:=\identity_L\otimes f$.
\end{lemma}
% segment-id: o014.aljabr2.chapter7.galois-descent-h1.pf001
\begin{proof}
	Pembaca dapat langsung memeriksa bahwa ${}^\sigma h$ memang pemetaan
	linear-$L$. Untuk kemulusan, pilih basis bagi $V$ dan $W$, lalu
	identifikasikan unsur-unsur $\Hom_L(V_L,W_L)$ dengan matriks. Grup
	$\Gal(L|K)$ kemudian beraksi dengan cara baku pada setiap entri matriks.
	Karena $h$ hanya melibatkan sejumlah hingga entri matriks, semuanya
	termuat dalam suatu subperluasan hingga $E|K$ dari $L|K$. Oleh karena itu,
	subgrup buka $\Gal(L|E)$ memfiksasi $h$. Pernyataan lain tentang aksi
	$(\sigma,h)\mapsto{}^\sigma h$ mudah diperiksa.

	Kesamaan
	$\Hom_L(V_L,W_L)^{\Gal(L|K)}=\Hom_K(V,W)$ merupakan akibat langsung
	penurunan Galois, dan juga mudah dilihat langsung dengan matriks.
\end{proof}

% segment-id: o014.aljabr2.chapter7.galois-descent-h1.p011
Dengan mengambil kasus khusus $V=W$, kita melihat bahwa $\Gal(L|K)$ beraksi
pada $\GL(V_L)$ melalui $(\sigma,g)\mapsto{}^\sigma g$, dan
$\GL(V_L)^{\Gal(L|K)}=\GL(V)$.

% segment-id: o014.aljabr2.chapter7.galois-descent-h1.lm002
\begin{lemma}
	Untuk data $\mathbf{t}$ berbentuk \eqref{eqn:tensor-data}, aksi mulus
	$\Gal(L|K)$ pada $\GL(V_L)$ mempertahankan $G(\mathbf{t}_L)$. Dengan
	demikian, $G(\mathbf{t}_L)$ menjadi grup-$\Gal(L|K)$ mulus dalam
	pengertian Catatan~\ref{rem:pro-finite-coh-nonabelian}; subgrup invariannya
	memenuhi
	% segment-id: o014.aljabr2.chapter7.galois-descent-h1.q010
	\[
	G(\mathbf{t}_L)^{\Gal(L|K)}=G(\mathbf{t}).
	\]
\end{lemma}
% segment-id: o014.aljabr2.chapter7.galois-descent-h1.pf002
\begin{proof}
	Misalkan $g\in G(\mathbf{t}_L)$ dan $\sigma\in\Gal(L|K)$. Maka
	% segment-id: o014.aljabr2.chapter7.galois-descent-h1.q011
	\[
	({}^\sigma g)\mathbf{t}_L
	=({}^\sigma g)(\sigma\mathbf{t}_L)
	=\sigma g\sigma^{-1}(\sigma\mathbf{t}_L)
	=\sigma(g\mathbf{t}_L)
	=\sigma(\mathbf{t}_L)
	=\mathbf{t}_L.
	\]
	Jadi, aksi $\Gal(L|K)$ mempertahankan $G(\mathbf{t}_L)$. Pernyataan
	lainnya langsung diperoleh.
\end{proof}

% segment-id: o014.aljabr2.chapter7.galois-descent-h1.p012
Dalam praktik, klasifikasi data \eqref{eqn:tensor-data} biasanya menjadi
lebih sederhana setelah memperluas medan ke medan $L$ yang cukup besar.
Sebagai contoh, kelas isomorfisme bentuk kuadratik takdegenerat di atas medan
tertutup aljabar sepenuhnya ditentukan oleh dimensinya, sedangkan situasinya
lebih rumit di atas medan umum seperti $\R$, $\Q_p$, atau $\Q$.

% segment-id: o014.aljabr2.chapter7.galois-descent-h1.p013
Syarat seperti ketakdegeneratan atau asosiativitas perkalian yang dikenakan
	dalam berbagai masalah klasifikasi di atas biasanya dapat diperiksa setelah
	sembarang perluasan medan. Karena itu, masalahnya terbagi menjadi dua bagian:
% segment-id: o014.aljabr2.chapter7.galois-descent-h1.l004
\begin{enumerate}
	\item Di atas medan $L$ yang cukup besar, klasifikasikan data berbentuk
	\eqref{eqn:tensor-data} yang memenuhi syarat yang diinginkan; pilihan
	khasnya adalah mengambil $L$ sebagai medan tertutup separabel.
	\item Pilih data $\mathbf{t}_0$ yang didefinisikan di atas $K$, ambil
	perluasan Galois $L|K$ yang cukup besar, lalu pelajari data $\mathbf{t}$
	yang memenuhi $\mathbf{t}_L\simeq\mathbf{t}_{0,L}$.
\end{enumerate}
Bagian pertama dapat disebut geometris, sedangkan bagian kedua aritmetis.
Bagian ini berfokus pada bagian aritmetis, yang memotivasi definisi berikut.

% segment-id: o014.aljabr2.chapter7.galois-descent-h1.d003
\begin{definition}
	Misalkan $\mathbf{t}_0$ data berbentuk \eqref{eqn:tensor-data} yang
	didefinisikan di atas medan $K$, dan misalkan $L|K$ suatu perluasan
	Galois. Tetapkan
	% segment-id: o014.aljabr2.chapter7.galois-descent-h1.q012
	\[
	\mathcal{T}(\mathbf{t}_0):=
	\left\{
	\begin{array}{r|l}
		\mathbf{t}:\text{ data berbentuk \eqref{eqn:tensor-data}}
		& \text{terdapat isomorfisme }h:\mathbf{t}_{0,L}\rightiso\mathbf{t}_L \\
		\text{yang didefinisikan di atas medan $K$}&
	\end{array}
	\right\}.
	\]
	Tuliskan pula
	$\mathscr{T}(\mathbf{t}_0):=\mathcal{T}(\mathbf{t}_0)/\simeq$.
\end{definition}

% segment-id: o014.aljabr2.chapter7.galois-descent-h1.p014
Definisi ini hanya mensyaratkan keberadaan isomorfisme $h$, tanpa menetapkan
pilihan tertentu; setiap dua pilihan berbeda melalui perkalian kanan dengan
suatu unsur $G(\mathbf{t}_{0,L})$. Tuliskan
% segment-id: o014.aljabr2.chapter7.galois-descent-h1.q013
\[
\Gamma:=\Gal(L|K),\qquad\text{dilengkapi topologi Krull}.
\]

% segment-id: o014.aljabr2.chapter7.galois-descent-h1.p015
Jika $\sigma\in\Gamma$ dan
$h:\mathbf{t}_{0,L}\rightiso\mathbf{t}_L$ (yakni
$h(\mathbf{t}_{0,L})=\mathbf{t}_L$), maka
% segment-id: o014.aljabr2.chapter7.galois-descent-h1.q014
\[
{}^\sigma h(\mathbf{t}_{0,L})
={}^\sigma h(\sigma\mathbf{t}_{0,L})
=\sigma(h(\mathbf{t}_{0,L}))
=\sigma(\mathbf{t}_L)
=\mathbf{t}_L.
\]
Karena itu, terdapat pemetaan tunggal
$c:\Gamma\to G(\mathbf{t}_{0,L})$ sedemikian sehingga
% segment-id: o014.aljabr2.chapter7.galois-descent-h1.q015
\[
{}^\sigma h=hc(\sigma),\qquad \sigma\in\Gal(L|K).
\]

% segment-id: o014.aljabr2.chapter7.galois-descent-h1.p016
Dari
${}^{\sigma\tau}h={}^\sigma(hc(\tau))={}^\sigma h\cdot
{}^\sigma c(\tau)=h\cdot c(\sigma)\cdot{}^\sigma c(\tau)$, diperoleh
% segment-id: o014.aljabr2.chapter7.galois-descent-h1.q016
\[
c(\sigma\tau)=c(\sigma)\cdot{}^\sigma c(\tau),
\qquad \sigma,\tau\in\Gamma.
\]

% segment-id: o014.aljabr2.chapter7.galois-descent-h1.p017
Tanpa mengurangi keumuman, andaikan $\mathbf{t}_0$ dan $\mathbf{t}$
masing-masing direalisasikan pada ruang $V$ dan $W$, sehingga
$h\in\Hom_L(V_L,W_L)$. Karena aksi $\Gamma$ mulus, terdapat subgrup buka
$\Gamma_0\subset\Gamma$ yang memfiksasi $h$; akibatnya, $c$ terfaktorkan
melalui $\Gamma/\Gamma_0$.

% segment-id: o014.aljabr2.chapter7.galois-descent-h1.p018
Jika pilihan $h$ diubah menjadi $h'=ha$, dengan
$a\in G(\mathbf{t}_{0,L})$, maka $c'$ yang bersesuaian memenuhi
% segment-id: o014.aljabr2.chapter7.galois-descent-h1.q017
\[
c'(\sigma)=(h')^{-1}\cdot{}^\sigma h'
=a^{-1}\cdot h^{-1}\cdot{}^\sigma h\cdot{}^\sigma a
=a^{-1}\cdot c(\sigma)\cdot{}^\sigma a.
\]

% segment-id: o014.aljabr2.chapter7.galois-descent-h1.p019
Dengan membandingkan hasil di atas satu demi satu dengan
Definisi~\ref{def:nonabelian-H} tentang $\Hm^1$ nonabelian, serta menggunakan
Catatan~\ref{rem:pro-finite-coh-nonabelian} untuk versi profinit, kita
memperoleh pemetaan
% segment-id: o014.aljabr2.chapter7.galois-descent-h1.q018
\begin{equation}\label{eqn:Galois-H1-prep}
	\begin{aligned}
		\mathcal{T}(\mathbf{t}_0)&\longrightarrow
		\Hm^1\left(\Gamma,G(\mathbf{t}_{0,L})\right),\\
		\mathbf{t}&\longmapsto[c]:=\text{kelas ekuivalensi }c.
	\end{aligned}
\end{equation}
Pemetaan ini kanonik dan tidak bergantung pada pilihan data bantu apa pun.

% segment-id: o014.aljabr2.chapter7.galois-descent-h1.p020
Ingat pula bahwa
$\Hm^1\left(\Gamma,G(\mathbf{t}_{0,L})\right)$ pada umumnya bukan grup,
tetapi merupakan himpunan bertitik yang titik pangkalnya diwakili oleh
pemetaan konstan $\Gamma\to G(\mathbf{t}_{0,L})$ dengan nilai automorfisme
identitas $1$. Di sisi lain, $\mathcal{T}(\mathbf{t}_0)$ juga mempunyai titik
pangkal yang jelas, yaitu $\mathbf{t}_0$, demikian pula himpunan hasil bagi
$\mathscr{T}(\mathbf{t}_0)$. Pemetaan \eqref{eqn:Galois-H1-prep} jelas
mempertahankan titik pangkal, sebagaimana tampak dengan mengambil
$h=\identity$.

% segment-id: o014.aljabr2.chapter7.galois-descent-h1.t001
\begin{theorem}\label{prop:Galois-H1}
	Untuk data $\mathbf{t}_0$ dan perluasan Galois $L|K$ seperti di atas,
	dengan $\Gamma:=\Gal(L|K)$, pemetaan \eqref{eqn:Galois-H1-prep}
	menginduksi bijeksi yang mempertahankan titik pangkal
	% segment-id: o014.aljabr2.chapter7.galois-descent-h1.q019
	\[
	\mathscr{T}(\mathbf{t}_0)\;\xrightarrow{1:1}\;
	\Hm^1\left(\Gamma,G(\mathbf{t}_{0,L})\right).
	\]
\end{theorem}
% segment-id: o014.aljabr2.chapter7.galois-descent-h1.pf003
\begin{proof}
	Pertama, kita tunjukkan bahwa \eqref{eqn:Galois-H1-prep} terfaktorkan melalui
	$\mathscr{T}(\mathbf{t}_0)$. Misalkan
	$f:\mathbf{t}\rightiso\mathbf{t}'$, sehingga
	$f_L:\mathbf{t}_L\rightiso\mathbf{t}'_L$. Dari
	$h:\mathbf{t}_{0,L}\rightiso\mathbf{t}_L$ diperoleh
	$f_Lh:\mathbf{t}_{0,L}\rightiso\mathbf{t}'_L$. Karena
	${}^\sigma f_L=f_L$, berlaku
	$(f_Lh)^{-1}\cdot{}^\sigma(f_Lh)=h^{-1}\cdot{}^\sigma h$, sehingga $c$
	yang bersesuaian tidak berubah.

	Selanjutnya, kita tunjukkan bahwa
	$\mathscr{T}(\mathbf{t}_0)\to
	\Hm^1\left(\Gamma,G(\mathbf{t}_{0,L})\right)$ injektif. Misalkan diberikan
	data $\mathbf{t}$, $\mathbf{t}'$, dan isomorfisme
	% segment-id: o014.aljabr2.chapter7.galois-descent-h1.q020
	\[
	\mathbf{t}'_L\xleftarrow[\sim]{h'}\mathbf{t}_{0,L}
	\xrightarrow[\sim]{h}\mathbf{t}_L,
	\]
	serta terdapat $a\in G(\mathbf{t}_{0,L})$ yang memenuhi
	% segment-id: o014.aljabr2.chapter7.galois-descent-h1.q021
	\[
	(h')^{-1}\cdot{}^\sigma h'
	=a^{-1}\cdot h^{-1}\cdot{}^\sigma h\cdot{}^\sigma a,
	\qquad\sigma\in\Gamma.
	\]
	Karena itu,
	% segment-id: o014.aljabr2.chapter7.galois-descent-h1.q022
	\[
	ha(h')^{-1}
	={}^\sigma h\cdot{}^\sigma a\cdot{}^\sigma(h')^{-1},
	\qquad\sigma\in\Gamma.
	\]
	Dengan kata lain,
	$ha(h')^{-1}:\mathbf{t}'_L\rightiso\mathbf{t}_L$ invarian terhadap aksi
	$\Gamma$, sehingga turun menjadi isomorfisme
	$\mathbf{t}'\rightiso\mathbf{t}$.

	Terakhir, kita tunjukkan bahwa
	$\mathscr{T}(\mathbf{t}_0)\to
	\Hm^1\left(\Gamma,G(\mathbf{t}_{0,L})\right)$ surjektif. Ambil
	$c\in Z^1(\Gamma,G(\mathbf{t}_{0,L}))$. Karena
	$c(\sigma)\in\GL(V_L)$, kita dapat menggunakannya untuk mendefinisikan
	aksi $\Gamma$ baru, yang dilambangkan dengan $\odot$, pada $V_L$:
	% segment-id: o014.aljabr2.chapter7.galois-descent-h1.q023
	\[
	\sigma\odot v:=c(\sigma)(\sigma v),
	\qquad\sigma\in\Gamma,\quad v\in V_L.
	\]

	Dari $c(\identity)=1$ dan
	% segment-id: o014.aljabr2.chapter7.galois-descent-h1.q024
	\[
	\sigma\odot(\tau\odot v)
	=c(\sigma)\sigma(c(\tau)(\tau v))
	=c(\sigma)\bigl({}^\sigma c(\tau)(\sigma\tau v)\bigr)
	=(\sigma\tau)\odot v,
	\]
	terlihat bahwa $\odot$ memang aksi $\Gamma$, sedangkan kemulusan $c$
	memastikan bahwa $\odot$ tetap merupakan aksi mulus dan semilinear.
	Definisikan ruang vektor-$K$
	$W:=(V_L)^{\Gamma,\;\odot\text{-aksi}}$. Berdasarkan penurunan Galois,
	kita memperoleh isomorfisme ruang vektor-$L$ yang mempertahankan aksi
	semilinear $\Gamma$,
	% segment-id: o014.aljabr2.chapter7.galois-descent-h1.q025
	\begin{gather*}
		h:(V_L,\;\odot\text{-aksi})\rightiso
		(W_L,\;\text{aksi alami}),\\
		w\in W\implies h(w)=1\otimes w.
	\end{gather*}

	Di sisi lain, $\Gamma$ juga beraksi secara semilinear melalui $\odot$ pada
	setiap $T^{p,q}(V_L)$ (lihat pembahasan setelah Definisi~\ref{def:tL}),
	dan karena $c(\sigma)\in G(\mathbf{t}_{0,L})$, berlaku
	% segment-id: o014.aljabr2.chapter7.galois-descent-h1.q026
	\[
	\sigma\odot\mathbf{t}_{0,L}
	=c(\sigma)(\sigma\mathbf{t}_{0,L})
	=c(\sigma)\mathbf{t}_{0,L}
	=\mathbf{t}_{0,L}.
	\]
	Karena itu, $\mathbf{t}_{0,L}$ turun relatif terhadap $\odot$ menjadi data
	$\mathbf{t}$ pada $W$, yang dicirikan oleh
	$h(\mathbf{t}_{0,L})=\mathbf{t}_L$.

	Dengan demikian, $\mathbf{t}\in\mathcal{T}(\mathbf{t}_0)$. Isomorfisme
	$h$ yang bersesuaian memenuhi
	$\sigma(h(v))=h(\sigma\odot v)=(hc(\sigma)\sigma)(v)$, yakni identitas
	$\sigma h=hc(\sigma)\sigma$. Menuliskannya kembali sebagai
	$h^{-1}\cdot{}^\sigma h=c(\sigma)$ memperlihatkan bahwa $\mathbf{t}$
	dipetakan ke $[c]$ oleh \eqref{eqn:Galois-H1-prep}. Ini membuktikan
	pernyataan tersebut.
\end{proof}

% segment-id: o014.aljabr2.chapter7.galois-descent-h1.p021
Penerapan yang paling sederhana dan langsung adalah sebagai berikut. Hasil
ini merupakan akibat dari fakta umum bahwa setiap ruang vektor berdimensi hingga
mempunyai basis, sehingga dimensi menentukan kelas isomorfisme ruang vektor
berdimensi hingga.

% segment-id: o014.aljabr2.chapter7.galois-descent-h1.t002
\begin{theorem}[Teorema Hilbert 90: versi $\GL$]\label{prop:Hilbert-90-GL}
	\index{Hilbert 90}
	Misalkan $V$ ruang vektor-$K$ berdimensi hingga dan $L|K$ perluasan
	Galois. Dengan demikian, $\Gamma:=\Gal(L|K)$ beraksi secara mulus pada
	grup $\GL(V_L)$. Maka
	% segment-id: o014.aljabr2.chapter7.galois-descent-h1.q027
	\[
	\Hm^1(\Gamma,\GL(V_L))=1.
	\]
\end{theorem}
% segment-id: o014.aljabr2.chapter7.galois-descent-h1.pf004
\begin{proof}
	Terapkan Teorema~\ref{prop:Galois-H1}, tetapi ambil data $\mathbf{t}_0$
	dengan $I=\emptyset$. Dengan kata lain, data tersebut hanya menentukan
	suatu ruang vektor-$K$ berdimensi hingga $V$, tanpa tensor apa pun. Dalam
	kasus ini, unsur-unsur $\mathscr{T}(\mathbf{t}_0)$ berkorespondensi
	satu-satu dengan kelas isomorfisme ruang vektor-$K$ $W$ yang memenuhi
	$W_L\simeq V_L$, atau ekuivalen dengan $\dim W=\dim V$. Jelas hanya ada
	satu kelas isomorfisme demikian.
\end{proof}

% segment-id: o014.aljabr2.chapter7.galois-descent-h1.p022
Serupa dengan itu, di atas medan berkarakteristik $\neq2$, fakta umum bahwa
setiap ruang simplektik $(V,\omega)$ mempunyai basis simplektik mengarah ke
hasil berikut.

% segment-id: o014.aljabr2.chapter7.galois-descent-h1.t003
\begin{theorem}
	Misalkan $\mathrm{char}(K)\neq2$, $V$ ruang vektor-$K$, dan
	$\omega:V\times V\to K$ bentuk bilinear antisimetris takdegenerat (juga
	disebut \emph{bentuk simplektik}, sedangkan data $(V,\omega)$ disebut
	\emph{ruang simplektik}). Grup simplektik yang bersesuaian didefinisikan
	sebagai
	% segment-id: o014.aljabr2.chapter7.galois-descent-h1.q028
	\begin{align*}
		\Sp(V)&=\Sp(V,\omega)\\
		&:=\{g\in\GL(V):\forall v_1,v_2\in V,\;
		\omega(gv_1,gv_2)=\omega(v_1,v_2)\}.
	\end{align*}

	Misalkan $L|K$ suatu perluasan Galois. Definisikan secara serupa grup
	$\Sp(V_L)$ beserta aksi $\Gamma:=\Gal(L|K)$ padanya; aksi ini mulus. Maka
	% segment-id: o014.aljabr2.chapter7.galois-descent-h1.q029
	\[
	\Hm^1(\Gamma,\Sp(V_L))=1.
	\]
\end{theorem}
% segment-id: o014.aljabr2.chapter7.galois-descent-h1.pf005
\begin{proof}
	Data $\mathbf{t}_0$ yang ditinjau terdiri atas ruang $V$ dan tensor
	$\omega$ bertipe $(0,2)$. Sifat takdegenerat dan antisimetris suatu bentuk
	bilinear dapat diperiksa setelah sembarang perluasan medan. Oleh karena itu,
	unsur-unsur $\mathscr{T}(\mathbf{t}_0)$ berkorespondensi satu-satu dengan
	kelas isomorfisme ruang simplektik $(W,\eta)$ di atas $K$ yang memenuhi
	$(W_L,\eta_L)\simeq(V_L,\omega_L)$. Berdasarkan teorema struktur untuk
	ruang simplektik, yakni bahwa dimensinya menentukan kelas isomorfismenya,
	data $(W,\eta)$ tersebut membentuk satu kelas isomorfisme saja.
\end{proof}

% segment-id: o014.aljabr2.chapter7.galois-descent-h1.eg001
\begin{example}\label{eg:multiplicative-90-finite}
	Ambil $\dim V=1$ dalam Teorema~\ref{prop:Hilbert-90-GL}. Karena
	$\GL(V_L)\simeq L^\times$ dengan aksi $\Gamma$ baku, diperoleh
	% segment-id: o014.aljabr2.chapter7.galois-descent-h1.q030
	\[
	\Hm^1(\Gamma,L^\times)=0.
	\]
	Karena $L^\times$ abelian, notasi aditif digunakan dalam persamaan ini.
	Jika $L|K$ hingga, persamaan tersebut dapat ditulis menggunakan kohomologi
	Tate dari Definisi~\ref{def:TaHm} sebagai
	$\TaHm^1(\Gamma,L^\times)=0$.

	Jika selanjutnya $L|K$ disyaratkan siklik, periodisitas yang dinyatakan
	dalam Contoh~\ref{eg:TaHm-periodic} memberikan
	% segment-id: o014.aljabr2.chapter7.galois-descent-h1.q031
	\[
	\TaHm^{-1}(\Gamma,L^\times)
	\simeq\TaHm^1(\Gamma,L^\times)=0.
	\]
	Dengan demikian, kita memperoleh kembali secara alami Teorema Hilbert 90
	versi multiplikatif; lihat \cite[Teorema 9.6.9]{Li1}.
\end{example}

% segment-id: o014.aljabr2.chapter7.galois-descent-h1.p023
Secara umum, menghitung $\Hm^1$ secara langsung untuk mengklasifikasikan
objek matematika mungkin cukup sulit. Manfaatnya terletak pada sudut pandang
baru yang diberikannya, serta pada tersedianya perangkat kohomologis seperti
barisan eksak panjang (Teorema~\ref{prop:nonabelian-long}) bagi kajian masalah
klasifikasi.
